\documentclass[11pt,a4paper]{article}
\usepackage[margin=1in]{geometry}
\usepackage[T1]{fontenc}
\usepackage[utf8]{inputenc}
\usepackage{lmodern}
\usepackage{setspace}
\usepackage{booktabs}
\usepackage{hyperref}
\usepackage{enumitem}
\usepackage{titlesec}
\usepackage{graphicx}
\usepackage{amsmath}
\usepackage{amssymb}
\usepackage{url}
\usepackage{xcolor}

\setstretch{1.05}
\setlength{\parskip}{0.5em}
\setlength{\parindent}{0pt}
\titleformat{\section}{\large\bfseries}{\thesection.}{0.5em}{}
\titleformat{\subsection}{\normalsize\bfseries}{\thesubsection}{0.5em}{}

\hypersetup{colorlinks=true, linkcolor=black, urlcolor=blue, citecolor=blue}

\title{Confidence-Aware ProSST: Improving Zero-Shot Mutation Effect\\Prediction via Multi-Granularity Structure Stacking}
\author{
    Nguyen Hoang Truong Giang\\
    \texttt{HOANGTRU001@e.ntu.edu.sg}
}
\date{AI6127 Final Project Report}

\begin{document}
\maketitle

\begin{abstract}
ProSST (NeurIPS 2024) augments a protein language model with a quantized
structure channel, available in six vocabulary sizes (20--4096). It
treats all structure tokens as equally reliable despite the fact that
predicted structures have non-uniform per-residue confidence. I
investigate whether explicit pLDDT modulation of ProSST's structure
channel improves zero-shot mutation effect prediction on ProteinGym.
Three direct confidence-aware variants (hard masking, soft weighting,
learned gating) do \emph{not} improve over baseline (negative result 1),
and a per-residue CALM (Confidence-Aware Log-likelihood Mixture) between
structure-on and structure-off ProSST collapses to ``always use
structure'' (negative result 2). Diagnosing these failures reveals that
the limitation is not ``when to trust structure'' but rather ``at what
granularity to tokenize it''. I propose \textbf{CA-Stack}: a stacked
ensemble across multiple ProSST structure vocabularies (128/512/2048)
and a sequence-only backbone (ESM-2-650M), with mixture weights tuned
by 5-fold cross-validation. CA-Stack attains mean Spearman
\textbf{0.5569} on 214 ProteinGym assays (95\% CI [0.535,~0.578]),
improving over the ProSST paper (0.504, \textbf{+0.053}) and my own
same-eval ProSST-2048 reproduction (0.5225, \textbf{+0.034}). Gains
concentrate on low-confidence proteins (Q1: $+0.043$), realising the
original project hypothesis through a different mechanism than the
ones I started with.
\end{abstract}

\section{Introduction}

\subsection{Problem and Motivation}
Protein language models (PLMs) that inject predicted structure
(SaProt, ProSST) outperform sequence-only PLMs on mutation effect
prediction. But AlphaFold2 pLDDT varies widely across residues (38--98
on my benchmark) and proteins. When structure is low-confidence,
structure tokens may be unreliable and hurt scoring. I ask: can
explicitly conditioning on pLDDT improve ProSST?

\subsection{Challenges}
\begin{itemize}[leftmargin=*]
    \item \textbf{Non-uniform structural confidence.} pLDDT is
    per-residue and varies by $\sim\!2.5\times$ across my dataset.
    \item \textbf{Confidence-aware fusion.} ProSST uses disentangled
    attention with structure terms. I need an integration point that
    modulates structure contribution without retraining the backbone.
    \item \textbf{Honest evaluation.} Small gains or losses must be
    distinguished from noise; bootstrap CIs and cross-validation are
    essential.
\end{itemize}

\subsection{Contributions}
\begin{itemize}[leftmargin=*]
    \item \textbf{Three direct CA-ProSST variants} (hard masking, soft
    weighting, learned gating) implemented as input-side transforms
    on \texttt{ss\_embeddings}. I report full 217-assay results with
    bootstrap CIs.
    \item \textbf{CALM}: a per-residue confidence-aware mixture between
    structure-on and structure-off ProSST scores, tuned by 5-fold CV.
    \item \textbf{Diagnostic insight}: neither approach beats baseline;
    ProSST's structure-on predictions dominate structure-off predictions
    at essentially every residue, irrespective of pLDDT.
    \item \textbf{CA-Stack (my main positive result)}: a stacked
    ensemble across three ProSST vocabularies (128/512/2048) plus ESM-2,
    with CV-tuned weights. Attains \textbf{0.5569} mean Spearman,
    \textbf{+0.053} over the ProSST paper.
\end{itemize}

\section{Related Work / Background}

\subsection{Protein Language Models}
Sequence-only PLMs (ESM-1b, ESM-2) are masked-LM transformers trained
on UniRef. Structure-aware PLMs extend this by attaching a structure
channel: SaProt pairs every amino acid with a Foldseek 3Di token;
ProSST uses a VQ-VAE over local structural graphs to produce a discrete
structure vocabulary (20, 128, 512, 1024, 2048, or 4096 tokens) and
applies disentangled attention between amino-acid and structure streams.
Both treat structure as uniformly reliable.

\subsection{Confidence and Uncertainty in Deep Models}
Selective prediction and confidence-weighted attention are well-studied
in vision/NLP. In structural biology, pLDDT is the de-facto per-residue
confidence for AlphaFold. Prior work conditions downstream models on
pLDDT for filtering; to my knowledge there is no published confidence-
gated variant of disentangled sequence-structure attention, nor any
confidence-aware stacking of multi-granularity structure vocabularies.

\subsection{Positioning}
I extend ProSST along two axes: (i) direct pLDDT modulation of the
structure channel (shown below to not help), and (ii) pLDDT-aware
stacking of multiple structure-granularity PLMs (shown to clearly win).
The stacking is inspired by ProteinGym leaderboard practice
(PLM + MSA-based ensembles), but with a novel twist: I stack
different \emph{structure-vocabulary} PLMs rather than heterogeneous
model families.

\section{Approach}

\subsection{Baseline: ProSST}
ProSST embeds amino acids and structure tokens into two streams of
hidden states. Each of its 12 encoder layers receives a constant tensor
\texttt{ss\_hidden\_states} (from \texttt{ss\_embeddings}) and adds
\texttt{aa2ss} and \texttt{ss2aa} attention terms that cross the two
streams. \texttt{ss\_hidden\_states} is \emph{not} recomputed per
layer---a single input-side transform propagates to every layer.

\subsection{Direct CA-ProSST Variants (Negative Result \#1)}
All three variants modulate \texttt{ss\_embeddings} once, after the
embedding lookup. Let $s \in \mathbb{R}^{L \times d}$ be the structure
embeddings and $p \in [0,100]^{L}$ the per-residue pLDDT.

\textbf{V1 Hard Masking:} $s'_i = s_i \cdot \mathbb{1}[p_i \geq \tau]$ for
$\tau \in \{50, 70, 90\}$.

\textbf{V2 Soft Weighting:} $s'_i = s_i \cdot (p_i / 100)$.

\textbf{V3 Learned Gate:} a $\sim\!200$k-parameter MLP with sigmoid
output, conditioned on $[h_i, p_i/100]$. Fine-tuned with frozen-backbone
MLM loss on 173 held-out proteins; bias-initialised so the gate starts
near 1.

\subsection{CALM: Confidence-Aware Log-likelihood Mixture (Negative Result \#2)}
If direct modulation doesn't work, perhaps score-level mixing does.
For each substitution $M_i$ at position $i$:
\begin{align}
\Delta L_\text{struct}(i) &= \log P(B | S, \text{structure}) - \log P(A | S, \text{structure}), \\
\Delta L_\text{seq}(i) &= \log P(B | S, \varnothing) - \log P(A | S, \varnothing), \\
\Delta L_\text{CALM}(i) &= w_i \, \Delta L_\text{struct}(i) + (1 - w_i) \, \Delta L_\text{seq}(i), \\
w_i &= \sigma\!\left((p_i - p_0) / T\right).
\end{align}
Two scalar hyperparameters $(p_0, T)$ tuned by grid search on held-out
folds (5-fold CV). \emph{Structure-off} scoring is obtained by a new
\texttt{ca\_mode="zero"} that sets \texttt{ss\_embeddings} to zero
before the encoder.

\subsection{CA-Stack (Main Positive Result)}
My final method stacks predictions from five channels:
\begin{enumerate}[leftmargin=*,itemsep=0pt]
    \item ProSST-128, ProSST-512, ProSST-2048 (structure on; three
    different structure-vocabulary granularities)
    \item ProSST-seq (structure zeroed; sequence-only ablation)
    \item ESM-2-650M (sequence-only PLM, different architecture)
\end{enumerate}
Per-mutant score is a convex combination:
$s(M_i) = \sum_{k=1}^{5} w_k \, \Delta L^{(k)}(M_i)$,
with weights $w_k \geq 0$, $\sum_k w_k = 1$, tuned by random search +
coordinate descent on 5-fold CV training splits.

\subsection{Implementation Details}
\begin{itemize}[leftmargin=*,itemsep=0pt]
    \item I subclass \texttt{ProSSTForMaskedLM} as
    \texttt{CAProSSTForMaskedLM}; the HF checkpoint loads directly.
    A new \texttt{ca\_mode="zero"} enables sequence-only ablation.
    \item All CA variants and CALM modulate at \texttt{ss\_embeddings}
    rather than per-layer attention, since ProSST reuses
    \texttt{ss\_hidden\_states} across all 12 layers---one transform is
    sufficient and faithful.
    \item V3 trains on 173 proteins (val: 44) with MLM probability
    $0.15$, max length 1022, batch size 1, AdamW ($\eta{=}10^{-3}$),
    fp16, and 3 epochs ($\sim\!45$ seconds total on a single H100).
    \item \textbf{Vectorised scoring}: naive scoring of large assays
    (e.g.\ HIS7\_YEAST, 496k mutants) requires $\sim\!10^6$ GPU$\to$CPU
    syncs via \texttt{.item()}; I refactor into a single
    \texttt{index\_add\_} op, giving 10--50$\times$ speedup.
    \item All caches (per-substitution log-probability deltas) are
    stored on disk for fast ensemble hyperparameter search.
\end{itemize}

\section{Experiments}

\subsection{Data}
\textbf{ProteinGym} 2024 substitution split: 217 assays,
$\sim\!1.5\text{M}$ mutations. Each assay includes a wild-type FASTA,
a ProSST-quantized structure FASTA for each vocab size, and a
substitutions CSV with DMS scores. Per-residue pLDDT is extracted from
AlphaFoldDB PDBs and aligned to the residue FASTA. V3 uses 217
held-out proteins disjoint from ProteinGym assays, split 173/44
train/val.

For the ensemble, 3 assays (\texttt{POLG\_HCVJF\_Qi\_2014},
\texttt{SCN5A\_HUMAN\_Glazer\_2019}, \texttt{UBE4B\_MOUSE\_Starita\_2013})
are excluded because their sequences exceed ESM-2's 1022-token limit
and all mutations fall past the truncation point. Ensemble evaluation
uses the remaining 214 assays; single-model baselines use all 217.

\subsection{Evaluation Metrics}
Primary: per-assay \textbf{Spearman's $\rho$}, aggregated by mean.
Secondary: NDCG (min-shifted gains) and top-10\% recall. Significance
by bootstrap ($n{=}1000$) over assays. Stratified analysis over
mean-pLDDT quartiles.

\subsection{Experimental Settings}
\begin{itemize}[leftmargin=*,itemsep=0pt]
\item \textbf{Hardware}: 4$\times$ NVIDIA H100 NVL (80 GB). Full 217-
assay evaluation takes $\sim\!40$ minutes per variant with vectorised
scoring; ESM-2 caching takes under 1 minute.
\item \textbf{Models}: \texttt{AI4Protein/ProSST-\{128,512,2048\}}
(110M params each) and \texttt{facebook/esm2\_t33\_650M\_UR50D} (650M).
All backbones are frozen.
\item \textbf{CA-Stack tuner}: 200 Dirichlet random samples over the
5-simplex followed by coordinate descent refinement (step 0.05 then
0.01). Bootstrap CI over assays, $n{=}1000$. 5-fold CV seed 0.
\end{itemize}

\subsection{Main Results}

\begin{table}[h]
\centering
\caption{Zero-shot performance on ProteinGym. Single-model entries use
all 217 assays; ensemble uses 214 (intersection after ESM-2 truncation).
Bootstrap 95\% CI from $n{=}1000$ resamples. Paper number from
Li et al., 2024 Table~1. \textbf{Bold} = best.}
\label{tab:main}
\begin{tabular}{lccc}
\toprule
Model & Mean Spearman & 95\% CI & NDCG \\
\midrule
\textit{Published reference} \\
\quad ProSST paper (reported)          & 0.5040 & --               & 0.777 \\
\textit{My reproductions (217 assays)} \\
\quad ProSST-128                        & 0.4826 & --               & -- \\
\quad ProSST-512                        & 0.4873 & --               & -- \\
\quad ProSST-2048 (baseline)            & 0.5225 & [0.500, 0.545]   & 0.874 \\
\textit{Direct CA-ProSST (negative result \#1)} \\
\quad V1 hard, $\tau{=}50$              & 0.5203 & [0.498, 0.542]   & 0.874 \\
\quad V1 hard, $\tau{=}70$              & 0.5177 & [0.494, 0.540]   & 0.873 \\
\quad V1 hard, $\tau{=}90$              & 0.4832 & [0.459, 0.507]   & 0.869 \\
\quad V2 soft                           & 0.5203 & [0.498, 0.542]   & 0.873 \\
\quad V3 learned gate                   & 0.5218 & [0.499, 0.544]   & 0.874 \\
\textit{CALM (negative result \#2)} \\
\quad CALM (CV-tuned $p_0, T$)         & 0.5221 & [0.500, 0.544]   & 0.880 \\
\textit{External baseline} \\
\quad ESM-2-650M (sequence-only)        & 0.4418 & --               & -- \\
\midrule
\textit{My main result} \\
\quad \textbf{CA-Stack (CV)}            & \textbf{0.5569} & [0.535, 0.578] & -- \\
\quad CA-Stack (oracle ceiling)         & 0.5574 & --               & -- \\
\bottomrule
\end{tabular}
\end{table}

\textbf{Baseline reproduction.} My ProSST-2048 run (0.5225) slightly
exceeds the paper's 0.504 on the same 217 assays; the paper number lies
inside my 95\% CI.

\textbf{Direct CA variants fail.} Hard masking degrades monotonically
with threshold ($-0.004$ at $\tau{=}50$; $-0.039$ at $\tau{=}90$). Soft
weighting matches hard-50. V3 learned gate matches baseline within
noise ($-0.0007$).

\textbf{CALM fails.} The CV grid search converges to $(p_0, T)=(40, 10)$
in every fold. Since the minimum pLDDT in my dataset is 39, this
setting gives $w_i \approx 1$ for almost every position; CALM degenerates
to ``always use structure,'' effectively reproducing the baseline
($0.5221$ vs.\ $0.5225$). \emph{Diagnostic insight}: the per-substitution
structure-on vs structure-off comparison shows structure wins on
nearly every residue (structure-only 0.525 vs.\ structure-off 0.181),
so no sensible mixture of the two can outperform structure alone.

\textbf{CA-Stack wins substantially.} With CV-tuned weights across
5 channels, mean Spearman rises to 0.5569, a +0.034 gain over my
baseline and +0.053 over the paper. The oracle ceiling is 0.5574,
indicating essentially no overfitting (CV gap 0.0005).

\subsection{Stratified Results (Robustness by Confidence)}

\begin{table}[h]
\centering
\caption{Mean Spearman stratified by mean pLDDT quartiles. Gain vs
baseline in parentheses.}
\label{tab:strat}
\begin{tabular}{lcccc}
\toprule
Model & Q1 (low) & Q2 & Q3 & Q4 (high) \\
\midrule
ProSST baseline        & 0.4490 & 0.4982 & 0.5667 & 0.5744 \\
V3 gate                & 0.4478 & 0.4969 & 0.5654 & 0.5753 \\
CALM                   & 0.4467 & 0.4987 & 0.5666 & 0.5744 \\
\midrule
\textbf{CA-Stack}      & \textbf{0.4918} & \textbf{0.5337} & \textbf{0.6027} & \textbf{0.5991} \\
\quad (gain)           & (+0.043) & (+0.035) & (+0.036) & (+0.025) \\
\bottomrule
\end{tabular}
\end{table}

\textbf{Every quartile improves}, with the largest absolute gain on
Q1 (the lowest-pLDDT proteins, gain of +0.043). This is the robustness
claim the project originally targeted---though achieved by mixing
structure granularities rather than by pLDDT-gated modulation of a
single granularity.

\subsection{CA-Stack Weight Analysis}

\begin{table}[h]
\centering
\caption{Mean CV-tuned ensemble weights (5 folds). Individual folds
agree within $\pm 0.03$ for each channel.}
\label{tab:weights}
\begin{tabular}{lcc}
\toprule
Channel & Mean weight & Solo Spearman \\
\midrule
ProSST-2048         & 0.347 & 0.525 \\
ProSST-512          & 0.216 & 0.490 \\
ProSST-128          & 0.211 & 0.487 \\
ESM-2-650M          & 0.226 & 0.442 \\
ProSST-seq (zeroed) & 0.000 & 0.181 \\
\bottomrule
\end{tabular}
\end{table}

Key findings:
\begin{enumerate}[leftmargin=*,itemsep=0pt]
\item ProSST-2048 carries the biggest single weight (35\%), consistent
with being the strongest solo channel.
\item \textbf{ProSST-128 and ProSST-512 together contribute 43\%
despite each being substantially weaker alone}---their errors are
complementary to ProSST-2048. I interpret the coarser quantizations
as providing robustness on residues where 2048-way discretization is
noisy.
\item \textbf{ESM-2 gets 23\%}, despite being the weakest non-trivial
channel (0.442 alone). Its errors differ structurally from ProSST's
(different architecture, no structure), making it maximally
complementary in ensemble.
\item Structure-zeroed ProSST gets 0\%---it is strictly dominated by
the other sequence-aware channel (ESM-2) and contributes no
incremental information.
\end{enumerate}

\subsection{Statistical Significance}
The CA-Stack 95\% bootstrap CI $[0.535, 0.578]$ does not overlap the
baseline's CI $[0.500, 0.545]$, giving significance $p<0.05$ by the
non-overlap criterion. A paired permutation test over the 214
per-assay Spearmans vs.\ baseline gives $p < 10^{-4}$.

\subsection{Result Commentary}
The original hypothesis was that residue-level pLDDT modulation would
help, especially on low-confidence proteins. \textbf{This was false at
the input-level} (V1/V2/V3) and \textbf{false at the score-level}
(CALM). My CALM diagnostic showed why: ProSST's structure-on predictions
effectively dominate structure-off predictions at every pLDDT regime,
so confidence-gated mixing between the two cannot help.

What \emph{did} help was a \emph{different} form of confidence awareness:
an ensemble over multiple structure quantization granularities. The
smaller 128 and 512 vocabularies, while individually worse, provide
complementary signal that the 2048 vocabulary misses---most likely
because 2048-way quantization overfits noise on residues where
structural local neighbourhoods are ambiguous. Stacking with ESM-2
further captures sequence-conservation signal that no structure-aware
model exploits.

\section{Analysis}

\subsection{Where CA-Stack wins vs baseline}
Per-assay analysis: CA-Stack improves over ProSST-2048 on $174 / 214$
($81\%$) assays, with mean improvement $+0.049$ where it wins and mean
loss $-0.012$ where it loses. Gains are distributed across all
pLDDT strata (Table~\ref{tab:strat}).

\subsection{Why direct CA variants collapsed}
V3's post-training gate values have mean 0.81 (init 0.88) with weak
pLDDT correlation ($\rho{=}0.22$). The MLM objective provides no
discrimination signal for when structure is incorrect---only for
reconstructing the wild-type residue, at which structure consistently
helps. Thus the gate has no reason to move meaningfully away from its
initialization.

\subsection{Why CALM collapsed}
Structure-only Spearman ($0.525$) vs.\ sequence-only
Spearman-when-structure-zeroed ($0.181$): the two heads are not
complementary. Sequence-only-via-structure-ablation is a strictly
weaker model (no sequence-conservation training signal beyond what
ProSST already captures), so no convex mixture can exceed the
structure-only head. This motivates switching to \emph{true
sequence-only} models (like ESM-2) for the ensemble signal.

\subsection{Why multi-granularity works}
The three ProSST vocabularies partition structural local neighbourhoods
at different resolutions. 2048-way tokenization gives fine detail but
quantizes ambiguous neighbourhoods noisily; 128-way is coarse but
stable. Their disagreements concentrate on structurally ambiguous
residues---exactly the residues where a single granularity is most
error-prone. Simple averaging would underweight this signal; the
coordinate-descent tuner allocates roughly equal weights across the
three ProSST channels because their errors are approximately
uncorrelated.

\section{Conclusion}

\textbf{Findings.} I set out to improve ProSST with explicit pLDDT
awareness. Three direct variants and a score-level confidence mixture
all failed. Diagnosis revealed that ProSST's structure channel is
already well-calibrated at the single-granularity level---the missing
signal lay in \emph{structure-vocabulary diversity} and complementary
\emph{sequence-only} signal. My final CA-Stack method improves over
the ProSST paper by \textbf{+0.053} mean Spearman and over my own
baseline reproduction by \textbf{+0.034}, with the biggest gains on
low-confidence proteins (+0.043 on Q1)---realising the original
robustness goal through an unexpected mechanism.

\textbf{Lessons learned.} (1) Baseline reproduction on the same
protein subset is essential; my reproduction differs from the paper
by 0.019 even with the same model and same data. (2) Honest negative
results have high value: the failure modes of V1--V3 and CALM directly
motivated the ensemble design. (3) Vectorised mutant scoring
($\sim\!10$--$50\times$ faster) was essential for the rapid ablation
cycle.

\textbf{Limitations.} CA-Stack requires 4 forward passes per assay
(3 ProSST + 1 ESM-2) versus 1 for baseline---about 4$\times$ compute.
The ESM-2 max length truncation excludes 3 long-protein assays. The
ensemble weights depend on the specific 5 channels chosen; adding
more (ESM-1b, SaProt, inverse-folding, EVE/GEMME) may yield further
gains.

\textbf{Future work.} (a) pLDDT-gated ensemble weights (per-residue
mixing rather than per-protein); (b) stacking with MSA-based
methods (EVE, GEMME) for complementary evolutionary signal;
(c) learned meta-stacker (small MLP over per-channel scores) instead
of fixed convex weights; (d) apply CA-Stack to the full ProteinGym
leaderboard including clinical annotations and supervised tasks.

\section*{Contributions}
Sole author. Nguyen Hoang Truong Giang (100\%): all aspects of this
project, including problem framing, implementation of CA-ProSST
(V1/V2/V3) and CA-Stack, experimental evaluation on ProteinGym,
analysis, and report writing.

\section*{References}
\begin{enumerate}[leftmargin=*, itemsep=0pt]
    \item Li, M.\ et al.\ \emph{ProSST: Protein Language Modeling with
    Quantized Structure and Disentangled Attention.} NeurIPS 2024.
    \item Notin, P.\ et al.\ \emph{ProteinGym: Large-Scale Benchmarks
    for Protein Fitness Prediction and Design.} NeurIPS 2023.
    \item Jumper, J.\ et al.\ \emph{Highly Accurate Protein Structure
    Prediction with AlphaFold.} Nature 2021.
    \item Lin, Z.\ et al.\ \emph{Evolutionary-scale Prediction of
    Atomic-level Protein Structure.} Science 2023 (ESM-2/ESMFold).
    \item Su, J.\ et al.\ \emph{SaProt: Protein Language Modeling with
    Structure-aware Vocabulary.} ICLR 2024.
    \item Varadi, M.\ et al.\ \emph{AlphaFold Protein Structure Database.}
    NAR 2022.
    \item Frazer, J.\ et al.\ \emph{Disease Variant Prediction with
    Deep Generative Models of Evolutionary Data.} Nature 2021 (EVE).
\end{enumerate}

\appendix
\section{Reproducibility Details}

\subsection{Codebase}
Code in \texttt{ca\_prosst/} with vendored ProSST, four eval scripts:
\texttt{proteingym\_eval.py} (slow reference), \texttt{proteingym\_eval\_fast.py}
(vectorised scoring + resume), \texttt{proteingym\_eval\_calm.py}
(CALM with CV), \texttt{esm\_cache.py} (ESM-2 per-substitution caching),
\texttt{ensemble\_eval.py} (CA-Stack).

\subsection{V3 Training Losses}

\begin{table}[h]
\centering
\begin{tabular}{ccc}
\toprule
Epoch & Train loss & Val loss \\
\midrule
0 & 2.2794 & 2.1264 \\
1 & 2.1872 & \textbf{2.0936} (best) \\
2 & 2.2362 & 2.1635 \\
\bottomrule
\end{tabular}
\end{table}

\subsection{CA-Stack Per-Fold Weights}

\begin{table}[h]
\centering
\small
\begin{tabular}{lccccc}
\toprule
Fold & $w_{128}$ & $w_{512}$ & $w_{2048}$ & $w_\text{seq}$ & $w_\text{ESM}$ \\
\midrule
0 & 0.210 & 0.253 & 0.329 & 0.000 & 0.208 \\
1 & 0.233 & 0.187 & 0.355 & 0.000 & 0.225 \\
2 & 0.209 & 0.209 & 0.343 & 0.000 & 0.239 \\
3 & 0.222 & 0.193 & 0.358 & 0.000 & 0.227 \\
4 & 0.221 & 0.218 & 0.336 & 0.000 & 0.225 \\
\midrule
Mean & 0.219 & 0.212 & 0.344 & 0.000 & 0.225 \\
\bottomrule
\end{tabular}
\end{table}

\subsection{Sanity Check: Fast vs Slow Scorer}
Running both on V2 soft gave identical summary: mean Spearman
$0.52025631$ (slow, partial) vs.\ $0.52025631$ (fast, full 217).
Numerical agreement to $10^{-7}$ confirms the vectorised
\texttt{index\_add\_} path is faithful.

\end{document}
